\documentclass[12pt,a4paper]{report}

% Packages
\usepackage[utf8]{inputenc}
\usepackage[T1]{fontenc}
\usepackage[lining]{ebgaramond}
\usepackage[cmintegrals,cmbraces]{newtxmath}
\usepackage{microtype}
\usepackage{geometry}
\usepackage{tabularx}
\usepackage{graphicx}
\usepackage{longtable}
\usepackage{booktabs}
\usepackage{array}
\usepackage{listings}
\usepackage{xcolor}
\usepackage{hyperref}
\usepackage{amsmath}
\usepackage{pdflscape}
\usepackage{float}
\usepackage{enumitem}
\usepackage{makecell}
\graphicspath{{./}}

% Page layout
\geometry{
    left=1.8cm,
    right=1.8cm,
    top=2.2cm,
    bottom=2.2cm
}

% Color definitions
\definecolor{crimson}{RGB}{153, 0, 0}
\definecolor{darkblue}{RGB}{0, 0, 153}

% Hyperlink settings
\hypersetup{
    colorlinks=true,
    linkcolor=crimson,
    urlcolor=crimson,
    citecolor=crimson
}

% Allow line breaks at slashes to prevent overfull hbox
\tolerance=2000
\emergencystretch=2em

% Chapter/section title formatting
\usepackage{titlesec}
\titleformat{\chapter}{\normalfont\huge\bfseries}{}{0pt}{}[\vspace{0.3em}]
\titlespacing*{\chapter}{0pt}{0pt}{14pt}
\titleformat{\section}{\normalfont\Large\bfseries}{\thesection}{1em}{}
\titleformat{\subsection}{\normalfont\large\bfseries}{\thesubsection}{1em}{}

\setlength{\parindent}{0pt}

\begin{document}

\pagestyle{plain}

% ─────────────────────────────────────────────
%  Cover Page
% ─────────────────────────────────────────────
\begin{titlepage}
\centering
\vspace*{3cm}

{\fontsize{28}{34}\selectfont\bfseries DSP in VLSI\par}
\vspace{0.6em}
{\fontsize{20}{26}\selectfont\bfseries Homework 4\par}
\vspace{0.4em}
{\fontsize{16}{22}\selectfont\bfseries Coordinate Rotation Digital Computer (CORDIC)\par}

\vspace{1.5cm}
\textcolor{crimson}{\rule{0.6\linewidth}{1.2pt}}
\vspace{1.5cm}

{\large\bfseries Student ID}\par
\vspace{0.4em}
{\large M11407439}\par

\vspace{0.8em}

{\large\bfseries Name}\par
\vspace{0.4em}
{\large MING-HONG, LIN}\par

\vspace{2cm}

{\large\bfseries Development Environment}\par
\vspace{0.8em}
\begin{tabular}{rl}
    \textbf{HDL}              & SystemVerilog \\[4pt]
    \textbf{Simulation Tools} & Synopsys VCS and Verdi \\[4pt]
    \textbf{Synthesis Tool}   & Synopsys DC Compiler \\[4pt]
    \textbf{Process}          & TSMC 16nm (ADFP) \\[4pt]
    \textbf{Verification}     & Python \\
\end{tabular}

\vfill

\textcolor{crimson}{\rule{0.6\linewidth}{0.6pt}}\par
\vspace{0.5em}
{\small Due Date: 2026/04/15}

\end{titlepage}

% ─────────────────────────────────────────────
%  Table of Contents
% ─────────────────────────────────────────────
\tableofcontents
\newpage

% ─────────────────────────────────────────────
\chapter{Step 1 — Realize CORDIC properties}
% ─────────────────────────────────────────────
\section{Result 1 - CORDIC scaling factor}

\hrule
\vspace{0.2cm}
Please draw the figure showing $S(N)$ \textit{versus} $N$ for $N = 1\ldots 30$. (5\%)
\vspace{0.2cm}
\hrule
\vspace{0.2cm}

% -------------------------
% USER WILL FILL IN HERE
% -------------------------

\newpage

% ─────────────────────────────────────────────
\chapter{Step 2 — Determine the word-length of fixed-point representation}
% ─────────────────────────────────────────────
\section{Result 2 - Word-length setting}

\hrule
\vspace{0.2cm}
Draw the figure of average absolute error versus fractional wordlength (10\%) to show how you determine the setting of word-length of the fractional part. Write down the integer word-length of $X(i)$ and $Y(i)$ that you use all the stages. Please explain it. (5\%)
\vspace{0.2cm}
\hrule
\vspace{0.2cm}

% -------------------------
% USER WILL FILL IN HERE
% -------------------------

\newpage

% ─────────────────────────────────────────────
\chapter{Step 3 — Determine the number of micro rotations for arctangent function}
% ─────────────────────────────────────────────
\section{Result 3 - Micro-rotations and elementary angles}

\hrule
\vspace{0.2cm}
Please draw a figure to denote the average phase errors of 10 quantized input pairs $(X, Y)$ versus different numbers of micro-rotations $S$ (10\%) and draw a figure to show the resulted phase errors of 10 quantized input pairs versus the word-length of quantized elementary angles (10\%). Explain your selection to satisfy the averaged error criterion of $2^{-9}$ radian. Also list a table of the elementary angles (both in floating-point representation and binary fixed-point representation). (5\%)
\vspace{0.2cm}
\hrule
\vspace{0.2cm}

% -------------------------
% USER WILL FILL IN HERE
% -------------------------

\newpage

% ─────────────────────────────────────────────
\chapter{Step 4 — Determine the number of micro rotations for magnitude function}
% ─────────────────────────────────────────────
\section{Result 4 - Number of micro-rotations for magnitude function}

\hrule
\vspace{0.2cm}
Please show how you decide the number of micro-rotations for the magnitude function with error tolerance of 0.1\%. (10\%)
\vspace{0.2cm}
\hrule
\vspace{0.2cm}

% -------------------------
% USER WILL FILL IN HERE
% -------------------------

\newpage

% ─────────────────────────────────────────────
\chapter{Step 5 — Determine the CSD of the scaling factor}
% ─────────────────────────────────────────────
\section{Result 5 - CSD representation and shift-and-add block}

\hrule
\vspace{0.2cm}
Write a program to show the setting of fractional wordlength of CSD versus error. Draw the figure. (10\%). Depict your design for the shift-and-add block according to your CSD representation. How many adders do you use? (10\%)
\vspace{0.2cm}
\hrule
\vspace{0.2cm}

% -------------------------
% USER WILL FILL IN HERE
% -------------------------

\newpage

% ─────────────────────────────────────────────
\chapter{Step 6 — Iterative architecture design}
% ─────────────────────────────────────────────
\section{Result 6.1 - Initial stage design}

\hrule
\vspace{0.2cm}
Depict your design of the initial stage. Mark the word-length in the block diagram. (10\%)
\vspace{0.2cm}
\hrule
\vspace{0.2cm}

% -------------------------
% USER WILL FILL IN HERE
% -------------------------

\newpage

\section{Result 6.2 - Iterative architecture behavior simulation}

\hrule
\vspace{0.2cm}
Implement your design with only DFFs inserted at the inputs and outputs. Show the timing diagram of behavior simulation for four inputs. (15\%) Compare the results with the arctangent function and draw the error versus index $m$ to show that your implementation meets the precision requirements of error less than $2^{-9}$. (10\%)
\vspace{0.2cm}
\hrule
\vspace{0.2cm}

% -------------------------
% USER WILL FILL IN HERE
% -------------------------

\newpage

% ─────────────────────────────────────────────
\chapter{Step 7 — Unfolding architecture design}
% ─────────────────────────────────────────────
\section{Result 7 - Unfolding architecture}

\hrule
\vspace{0.2cm}
Draw the block diagram of your $\frac{S}{2}$-unfolding architecture. (10\%). Show the timing diagram of behavior and post-synthesis simulation with correct setting of clock period and 10 inputs. (30\%) Verified the error between the post-synthesis results and the floating-point arctangent results. Draw the error versus index $m$. (10\%)
\vspace{0.2cm}
\hrule
\vspace{0.2cm}

% -------------------------
% USER WILL FILL IN HERE
% -------------------------

\newpage

% ─────────────────────────────────────────────
\chapter{Step 8 — Check area}
% ─────────────────────────────────────────────
\section{Result 8 - Area comparison}

\hrule
\vspace{0.2cm}
Report and compare the area before and after unfolding. (10\%)
\vspace{0.2cm}
\hrule
\vspace{0.2cm}

% -------------------------
% USER WILL FILL IN HERE
% -------------------------

\newpage

% ─────────────────────────────────────────────
\chapter{Step 9 — Hardware architecture design}
% ─────────────────────────────────────────────
\section{Result 9 - Architecture for magnitude function}

\hrule
\vspace{0.2cm}
Draw the shift-and-add module for the scaling factor (5\%). Connect the shift-and-add module to your micro-rotation module, and calculate the magnitude function. Show the timing diagram of behavior simulation and post-synthesis simulation with proper setting of clock period. (15\%). Verified the error between the post-synthesis results and the floating-point arctangent results. Draw the error versus index $m$. (10\%)
\vspace{0.2cm}
\hrule
\vspace{0.2cm}

% -------------------------
% USER WILL FILL IN HERE
% -------------------------

\end{document}
